\section{Introduction}
\label{sec:intro}

The United States has seen skyrocketing real estate and rental prices in recent decades in its major metropolitan areas, and these increases have outpaced median income growth \cite{joint_center_for_housing_studies_of_harvard_university_americas_2024}.
In particular, New York City has long been seen as the most expensive city in the country and has gotten even more expensive in recent years \cite{glaeser_why_2005-1}.
Many studies have highlighted the role of land-use regulation, and in particular discretionary land-use decision-making, in lowering housing supply growth and leading to rising prices \cite{fischel_homevoter_2009,hilber_origins_2013, paciorek_supply_2013, khan_decentralized_2021,kulka_how_2022}.

This paper studies whether increased local control affects housing supply using a sharp institutional change in New York City's land-use process. 
Until 1989, the Uniform Land Use Review Procedure ended with the Board of Estimate, a citywide body made up of the mayor, comptroller, city council president, and the five borough presidents. 
In \textit{Board of Estimate of City of New York v. Morris}, the U.S. Supreme Court held that the Board's structure violated one-person, one-vote because boroughs with very different populations received equal representation through their borough presidents \cite{board_of_estimate_of_city_of_new_york_v_morris_notitle_1989}. 
The 1989 Charter revision, induced by the \textit{Morris} decision, drastically changed that allocation of authority. 
After the Board was abolished, final review of many land-use actions shifted to the City Council \citep{lane_policy_1998}. 
The Council is a 51-member districted legislature. 
A land-use application is still reviewed by the community board, borough president, and City Planning Commission, but the Council now sits at the final legislative stage for many discretionary approvals. The mayor can veto some Council actions, but the practical politics of district-specific land use are dominated by the Council.

Although the Charter gives formal land-use authority to the Council as a whole, the Council generally follows the local member on district-specific land-use applications \cite{donaldson_charter_2025}.
In practice, this gives each councilmember substantial control over projects in the member's district. 
If homeowners are the local group most likely to oppose large new buildings \cite{fischel_homevoter_2009}, then the shift from borough-scale review to district-scale review should matter most in places where homeowners were especially concentrated relative to the rest of their borough.

I measure that likelihood to oppose new developments using 1990 homeownership in New York City's community districts (CDs), which is a proxy for constituencies who are more likely to oppose large new developments, particularly of rental or affordable housing \cite{khan_decentralized_2021}.
Community districts are useful because they are stable neighborhood units observed before and after the reform, unlike city council districts. There are also 59 of them relative to 51 city council districts, so they are very similar in their level of granularity. 
The main treatment measure I construct is a district's 1990 homeownership share minus its borough's 1990 homeownership share, standardized within borough, which allows me to measure variation in local opposition without conflating borough-level differences in homeownership and development activity.
A high-homeownership district in Manhattan is not politically or physically comparable to a high-homeownership district in Staten Island. 
The treatment instead asks whether, within the same borough, homeowner-heavy community districts experienced different post-reform housing trajectories than nearby districts with lower homeowner concentration.

I find suggestive evidence for my hypothesis that higher homeowner CDs saw lower flows of new residential development after the political transition occured. 
Raw trends show that after 2010, the high-treatment districts account for a much smaller share of large-building construction than the low-treatment districts. 
In 2010--2019, the high-treatment tercile accounts for 15 percent of borough-total 50+ unit construction, compared with 43 percent in the low-treatment tercile. In 2020--2025, the corresponding shares are 13 percent and 49 percent. 
I also use Brooklyn as a case study that shows the same pattern in a more concrete way: high-treatment Brooklyn districts produce fewer large-building units, while the small-building margin moves little or in the opposite direction.
